% --- Archon physics formalization source begin ---
% archon:physics
% archon:covers PhyXMiniProblems/problem_phyx_mini_0250.lean
% archon:source-report reports/phyx_mini/problem_phyx_mini_0250.source.json

\chapter{Physics problem phyx\_mini\_0250}
\label{ch:PhyXMiniProblems_problem_phyx_mini_0250}

\paragraph{Problem source.}
\#\# Physical scenario

A Michelson interferometer uses a helium-neon laser with wavelength \$\textbackslash{}lambda\_\{vac\} = 633\textbackslash{};nm\$. In one arm, the light passes through a \$4.00\textbackslash{};cm\$-thick glass cell. Initially the cell is evacuated, and the interferometer is adjusted so that the central spot is a bright fringe. The cell is then slowly filled to atmospheric pressure with a gas. As the cell fills, 43 bright-dark-bright fringe shifts are seen and counted.

\#\# Question

What is the index of refraction of the gas at this wavelength?

\#\# Answer choices

- A: A: 1.00034

- B: B: 1.00031

- C: C: 1.00042

- D: D: 1.00028

\#\# Figure caption (auxiliary; use the image as primary evidence)

The image is a diagram of an optical setup, likely an interferometer. Here is a breakdown of its components:

1. **Source**: The origin of the light, depicted on the left side.

2. **Beam Splitter**: Positioned in the path of the light from the source. It splits the light into two paths.

3. **Mirror \textbackslash{}(M\_1\textbackslash{})**: Located at the top of the image, it reflects one of the split beams back through the beam splitter.

4. **Mirror \textbackslash{}(M\_2\textbackslash{})**: Placed on the right side. It also reflects the light back and is shown with an arrow indicating its possible movement.

5. **Gas-filled Cell**: Placed in the path leading to \textbackslash{}(M\_2\textbackslash{}) and labeled as having a thickness \textbackslash{}(d\textbackslash{}). The light passes through this cell twice, as noted in the text.

6. **Path Lengths**: The paths have labeled distances \textbackslash{}(L\_1\textbackslash{}) and \textbackslash{}(L\_2\textbackslash{}).

7. **Eye**: An icon at the bottom indicating the observation point or detector for the combined light beams.

8. **Text**: Above the gas-filled cell, there is a description stating: "Gas-filled cell of thickness d. Light goes through this cell twice."

The diagram shows the relationships between these components, indicating how light is manipulated and directed through the setup.

\#\# Dataset metadata

Wave Properties; Physical Model Grounding Reasoning, Multi-Formula Reasoning

\paragraph{Recorded answer/context.}
A: A: 1.00034

\paragraph{Figure/image path.}
/root/proposal\_for\_physic/science-mango/phyx\_mini\_run/phyx\_data/test\_image/250.png

\paragraph{Formalization target.}
create a compiling Lean file with sorry bodies at `PhyXMiniProblems/problem\_phyx\_mini\_0250.lean`. The Lean declarations must preserve the physical quantities, dimensions or dimensional roles, figure labels, governing-law hypotheses, and final relation expressed by this problem.
Use Mathlib/Physlib names found through LeanExplore where available. If a physics API is missing, introduce faithful local abstractions rather than scalar placeholder aliases.

\begin{theorem}[Physics formalization target]
\label{thm:physics:phyx_mini_0250:target}
\lean{PhyXMiniProblems.ProblemPhyXMini0250.problem_phyx_mini_0250}
\uses{def:physics:phyx-mini-0250:phyxminiproblems-problemphyxmini0250-michelsongascellexperiment, def:physics:phyx-mini-0250:phyxminiproblems-problemphyxmini0250-hasstatedproblemreadouts, def:physics:phyx-mini-0250:phyxminiproblems-problemphyxmini0250-matchessourcefigure, def:physics:phyx-mini-0250:phyxminiproblems-problemphyxmini0250-hasphysicalparameters, def:physics:phyx-mini-0250:phyxminiproblems-problemphyxmini0250-hasvacuumreferenceindex, def:physics:phyx-mini-0250:phyxminiproblems-problemphyxmini0250-obeysroundtripopticalpathlaw, def:physics:phyx-mini-0250:phyxminiproblems-problemphyxmini0250-obeyscompletefringeshiftlaw, def:physics:phyx-mini-0250:phyxminiproblems-problemphyxmini0250-answerrefractiveindex, def:physics:phyx-mini-0250:phyxminiproblems-problemphyxmini0250-isnearesthundredthousandthreadout}
The assigned autoformalize agent should translate this physics problem into Lean declarations in the covered file, with theorem and lemma proof bodies written as `by sorry`.
\end{theorem}
\begin{proof}
This is an autoformalization task, not a proof task. Produce statements that can later be proved by the physics prover without weakening the physical model.
\end{proof}

% --- Archon declaration topology begin ---
\section{Declaration topology}
\begin{definition}[Length Quantity]
  \label{def:physics:phyx-mini-0250:phyxminiproblems-problemphyxmini0250-lengthquantity}
  \lean{PhyXMiniProblems.ProblemPhyXMini0250.LengthQuantity}
  A physical length represented independently of the chosen unit system
\end{definition}
\begin{proof}
  This is the declared model definition of Length Quantity.
\end{proof}
\begin{definition}[length Readout]
  \label{def:physics:phyx-mini-0250:phyxminiproblems-problemphyxmini0250-lengthreadout}
  \lean{PhyXMiniProblems.ProblemPhyXMini0250.lengthReadout}
  \uses{def:physics:phyx-mini-0250:phyxminiproblems-problemphyxmini0250-lengthquantity}
  Read a physical length as a real scalar in the selected length unit
\end{definition}
\begin{proof}
  This is the declared model definition of length Readout.
\end{proof}
\begin{definition}[length In Meters]
  \label{def:physics:phyx-mini-0250:phyxminiproblems-problemphyxmini0250-lengthinmeters}
  \lean{PhyXMiniProblems.ProblemPhyXMini0250.lengthInMeters}
  \uses{def:physics:phyx-mini-0250:phyxminiproblems-problemphyxmini0250-lengthquantity, def:physics:phyx-mini-0250:phyxminiproblems-problemphyxmini0250-lengthreadout}
  The metre readout used by the governing optical-path laws
\end{definition}
\begin{proof}
  This is the declared model definition of length In Meters.
\end{proof}
\begin{definition}[length In Centimeters]
  \label{def:physics:phyx-mini-0250:phyxminiproblems-problemphyxmini0250-lengthincentimeters}
  \lean{PhyXMiniProblems.ProblemPhyXMini0250.lengthInCentimeters}
  \uses{def:physics:phyx-mini-0250:phyxminiproblems-problemphyxmini0250-lengthquantity, def:physics:phyx-mini-0250:phyxminiproblems-problemphyxmini0250-lengthreadout}
  The centimetre readout used for the cell thickness d
\end{definition}
\begin{proof}
  This is the declared model definition of length In Centimeters.
\end{proof}
\begin{definition}[length In Nanometers]
  \label{def:physics:phyx-mini-0250:phyxminiproblems-problemphyxmini0250-lengthinnanometers}
  \lean{PhyXMiniProblems.ProblemPhyXMini0250.lengthInNanometers}
  \uses{def:physics:phyx-mini-0250:phyxminiproblems-problemphyxmini0250-lengthquantity, def:physics:phyx-mini-0250:phyxminiproblems-problemphyxmini0250-lengthreadout}
  The nanometre readout used for the helium-neon laser wavelength
\end{definition}
\begin{proof}
  This is the declared model definition of length In Nanometers.
\end{proof}
\begin{definition}[Laser Kind]
  \label{def:physics:phyx-mini-0250:phyxminiproblems-problemphyxmini0250-laserkind}
  \lean{PhyXMiniProblems.ProblemPhyXMini0250.LaserKind}
  The source type named in the problem statement
\end{definition}
\begin{proof}
  This is the declared model definition of Laser Kind.
\end{proof}
\begin{definition}[Cell State]
  \label{def:physics:phyx-mini-0250:phyxminiproblems-problemphyxmini0250-cellstate}
  \lean{PhyXMiniProblems.ProblemPhyXMini0250.CellState}
  The two cell states compared during the slow filling experiment
\end{definition}
\begin{proof}
  This is the declared model definition of Cell State.
\end{proof}
\begin{definition}[Filling Protocol]
  \label{def:physics:phyx-mini-0250:phyxminiproblems-problemphyxmini0250-fillingprotocol}
  \lean{PhyXMiniProblems.ProblemPhyXMini0250.FillingProtocol}
  The quasi-static preparation described by the prose
\end{definition}
\begin{proof}
  This is the declared model definition of Filling Protocol.
\end{proof}
\begin{definition}[Fringe Brightness]
  \label{def:physics:phyx-mini-0250:phyxminiproblems-problemphyxmini0250-fringebrightness}
  \lean{PhyXMiniProblems.ProblemPhyXMini0250.FringeBrightness}
  Brightness states of the central interference spot
\end{definition}
\begin{proof}
  This is the declared model definition of Fringe Brightness.
\end{proof}
\begin{definition}[Fringe Shift Pattern]
  \label{def:physics:phyx-mini-0250:phyxminiproblems-problemphyxmini0250-fringeshiftpattern}
  \lean{PhyXMiniProblems.ProblemPhyXMini0250.FringeShiftPattern}
  The complete fringe pattern counted once per observed shift
\end{definition}
\begin{proof}
  This is the declared model definition of Fringe Shift Pattern.
\end{proof}
\begin{definition}[Interferometer Arm]
  \label{def:physics:phyx-mini-0250:phyxminiproblems-problemphyxmini0250-interferometerarm}
  \lean{PhyXMiniProblems.ProblemPhyXMini0250.InterferometerArm}
  The two Michelson arms, carrying the figure labels L₁ and L₂
\end{definition}
\begin{proof}
  This is the declared model definition of Interferometer Arm.
\end{proof}
\begin{definition}[Figure Element]
  \label{def:physics:phyx-mini-0250:phyxminiproblems-problemphyxmini0250-figureelement}
  \lean{PhyXMiniProblems.ProblemPhyXMini0250.FigureElement}
  Optical components and the observer shown in the primary figure
\end{definition}
\begin{proof}
  This is the declared model definition of Figure Element.
\end{proof}
\begin{definition}[Beam Leg]
  \label{def:physics:phyx-mini-0250:phyxminiproblems-problemphyxmini0250-beamleg}
  \lean{PhyXMiniProblems.ProblemPhyXMini0250.BeamLeg}
  Directed beam legs drawn in the primary figure
\end{definition}
\begin{proof}
  This is the declared model definition of Beam Leg.
\end{proof}
\begin{definition}[Propagation Direction]
  \label{def:physics:phyx-mini-0250:phyxminiproblems-problemphyxmini0250-propagationdirection}
  \lean{PhyXMiniProblems.ProblemPhyXMini0250.PropagationDirection}
  Cardinal directions used to transcribe the beam arrows
\end{definition}
\begin{proof}
  This is the declared model definition of Propagation Direction.
\end{proof}
\begin{definition}[Michelson Gas Cell Experiment]
  \label{def:physics:phyx-mini-0250:phyxminiproblems-problemphyxmini0250-michelsongascellexperiment}
  \lean{PhyXMiniProblems.ProblemPhyXMini0250.MichelsonGasCellExperiment}
  \uses{def:physics:phyx-mini-0250:phyxminiproblems-problemphyxmini0250-lengthquantity, def:physics:phyx-mini-0250:phyxminiproblems-problemphyxmini0250-laserkind, def:physics:phyx-mini-0250:phyxminiproblems-problemphyxmini0250-cellstate, def:physics:phyx-mini-0250:phyxminiproblems-problemphyxmini0250-fillingprotocol, def:physics:phyx-mini-0250:phyxminiproblems-problemphyxmini0250-fringebrightness, def:physics:phyx-mini-0250:phyxminiproblems-problemphyxmini0250-fringeshiftpattern, def:physics:phyx-mini-0250:phyxminiproblems-problemphyxmini0250-interferometerarm, def:physics:phyx-mini-0250:phyxminiproblems-problemphyxmini0250-figureelement, def:physics:phyx-mini-0250:phyxminiproblems-problemphyxmini0250-beamleg, def:physics:phyx-mini-0250:phyxminiproblems-problemphyxmini0250-propagationdirection}
  The physical quantities and figure data of the interferometer. refractiveIndex is an unknown dimensionless material readout at each cell state. The field does not contain the requested numerical answer
\end{definition}
\begin{proof}
  This is the declared model definition of Michelson Gas Cell Experiment.
\end{proof}
\begin{definition}[Has Stated Problem Readouts]
  \label{def:physics:phyx-mini-0250:phyxminiproblems-problemphyxmini0250-hasstatedproblemreadouts}
  \lean{PhyXMiniProblems.ProblemPhyXMini0250.HasStatedProblemReadouts}
  \uses{def:physics:phyx-mini-0250:phyxminiproblems-problemphyxmini0250-lengthincentimeters, def:physics:phyx-mini-0250:phyxminiproblems-problemphyxmini0250-lengthinnanometers, def:physics:phyx-mini-0250:phyxminiproblems-problemphyxmini0250-michelsongascellexperiment}
  This def records the Has Stated Problem Readouts component of the problem-specific physical model
\end{definition}
\begin{proof}
  This is the declared model definition of Has Stated Problem Readouts.
\end{proof}
\begin{definition}[Matches Source Figure]
  \label{def:physics:phyx-mini-0250:phyxminiproblems-problemphyxmini0250-matchessourcefigure}
  \lean{PhyXMiniProblems.ProblemPhyXMini0250.MatchesSourceFigure}
  \uses{def:physics:phyx-mini-0250:phyxminiproblems-problemphyxmini0250-michelsongascellexperiment}
  This def records the Matches Source Figure component of the problem-specific physical model
\end{definition}
\begin{proof}
  This is the declared model definition of Matches Source Figure.
\end{proof}
\begin{definition}[Has Physical Parameters]
  \label{def:physics:phyx-mini-0250:phyxminiproblems-problemphyxmini0250-hasphysicalparameters}
  \lean{PhyXMiniProblems.ProblemPhyXMini0250.HasPhysicalParameters}
  \uses{def:physics:phyx-mini-0250:phyxminiproblems-problemphyxmini0250-lengthinmeters, def:physics:phyx-mini-0250:phyxminiproblems-problemphyxmini0250-michelsongascellexperiment}
  Positivity and monotonicity conditions for the physical experiment
\end{definition}
\begin{proof}
  This is the declared model definition of Has Physical Parameters.
\end{proof}
\begin{definition}[Has Vacuum Reference Index]
  \label{def:physics:phyx-mini-0250:phyxminiproblems-problemphyxmini0250-hasvacuumreferenceindex}
  \lean{PhyXMiniProblems.ProblemPhyXMini0250.HasVacuumReferenceIndex}
  \uses{def:physics:phyx-mini-0250:phyxminiproblems-problemphyxmini0250-michelsongascellexperiment}
  The evacuated reference state has vacuum refractive index one
\end{definition}
\begin{proof}
  This is the declared model definition of Has Vacuum Reference Index.
\end{proof}
\begin{definition}[Obeys Round Trip Optical Path Law]
  \label{def:physics:phyx-mini-0250:phyxminiproblems-problemphyxmini0250-obeysroundtripopticalpathlaw}
  \lean{PhyXMiniProblems.ProblemPhyXMini0250.ObeysRoundTripOpticalPathLaw}
  \uses{def:physics:phyx-mini-0250:phyxminiproblems-problemphyxmini0250-lengthinmeters, def:physics:phyx-mini-0250:phyxminiproblems-problemphyxmini0250-michelsongascellexperiment}
  This def records the Obeys Round Trip Optical Path Law component of the problem-specific physical model
\end{definition}
\begin{proof}
  This is the declared model definition of Obeys Round Trip Optical Path Law.
\end{proof}
\begin{definition}[Obeys Complete Fringe Shift Law]
  \label{def:physics:phyx-mini-0250:phyxminiproblems-problemphyxmini0250-obeyscompletefringeshiftlaw}
  \lean{PhyXMiniProblems.ProblemPhyXMini0250.ObeysCompleteFringeShiftLaw}
  \uses{def:physics:phyx-mini-0250:phyxminiproblems-problemphyxmini0250-lengthinmeters, def:physics:phyx-mini-0250:phyxminiproblems-problemphyxmini0250-michelsongascellexperiment}
  This def records the Obeys Complete Fringe Shift Law component of the problem-specific physical model
\end{definition}
\begin{proof}
  This is the declared model definition of Obeys Complete Fringe Shift Law.
\end{proof}
\begin{definition}[Answer Choice]
  \label{def:physics:phyx-mini-0250:phyxminiproblems-problemphyxmini0250-answerchoice}
  \lean{PhyXMiniProblems.ProblemPhyXMini0250.AnswerChoice}
  Labels of the four numerical choices printed with the problem
\end{definition}
\begin{proof}
  This is the declared model definition of Answer Choice.
\end{proof}
\begin{definition}[answer Refractive Index]
  \label{def:physics:phyx-mini-0250:phyxminiproblems-problemphyxmini0250-answerrefractiveindex}
  \lean{PhyXMiniProblems.ProblemPhyXMini0250.answerRefractiveIndex}
  \uses{def:physics:phyx-mini-0250:phyxminiproblems-problemphyxmini0250-answerchoice}
  The dimensionless refractive-index readout displayed by each choice
\end{definition}
\begin{proof}
  This is the declared model definition of answer Refractive Index.
\end{proof}
\begin{definition}[Is Nearest Hundred Thousandth Readout]
  \label{def:physics:phyx-mini-0250:phyxminiproblems-problemphyxmini0250-isnearesthundredthousandthreadout}
  \lean{PhyXMiniProblems.ProblemPhyXMini0250.IsNearestHundredThousandthReadout}
  Agreement with a five-decimal display by nearest-value rounding
\end{definition}
\begin{proof}
  This is the declared model definition of Is Nearest Hundred Thousandth Readout.
\end{proof}
\begin{lemma}[gas Refractive Index eq exact]
  \label{lem:physics:phyx-mini-0250:phyxminiproblems-problemphyxmini0250-gasrefractiveindex-eq-exact}
  \lean{PhyXMiniProblems.ProblemPhyXMini0250.gasRefractiveIndex_eq_exact}
  \uses{def:physics:phyx-mini-0250:phyxminiproblems-problemphyxmini0250-michelsongascellexperiment, def:physics:phyx-mini-0250:phyxminiproblems-problemphyxmini0250-hasstatedproblemreadouts, def:physics:phyx-mini-0250:phyxminiproblems-problemphyxmini0250-matchessourcefigure, def:physics:phyx-mini-0250:phyxminiproblems-problemphyxmini0250-hasphysicalparameters, def:physics:phyx-mini-0250:phyxminiproblems-problemphyxmini0250-hasvacuumreferenceindex, def:physics:phyx-mini-0250:phyxminiproblems-problemphyxmini0250-obeysroundtripopticalpathlaw, def:physics:phyx-mini-0250:phyxminiproblems-problemphyxmini0250-obeyscompletefringeshiftlaw}
  This lemma records the gas Refractive Index eq exact component of the problem-specific physical model
\end{lemma}
\begin{proof}
  The conclusion follows from the governing relations and figure readouts named in the statement of gas Refractive Index eq exact.
\end{proof}
% --- Archon declaration topology end ---
% --- Archon physics formalization source end ---
